% ============================================================
% Reporte: mapeo-memoria
% Generado a partir de reportes/comun/plantilla_base.tex
% (NO edites la plantilla: cada unidad trabaja en su main.tex)
%
% Reglas de secciones y extensiones: reportes/AGENTS.md
% Compilar SIEMPRE con:
%   ../../scripts/latexbuild main.tex      (desde esta carpeta)
%   → genera build/main.pdf
% ============================================================
\documentclass[llncs]{llncs}

\usepackage[utf8]{inputenc}
\usepackage[T1]{fontenc}

% ---- AVISO DE CODIFICACIÓN (xelatex + T1) ----
% Con la cadena del proyecto (xelatex con codificación T1) los caracteres
% UTF-8 de la banda 0x80-0xFF escritos como literal se remapean mal y
% salen otros glifos: « » § ª · aparecen como ń ż ğ ł ů.
% En su lugar usa las macros de LaTeX, que sí funcionan:
%   «texto»  ->  ``texto''          §6.2   ->  \S6.2
%   ª        ->  \textordfeminine   ·      ->  \textperiodcentered
% Sí funcionan como literal: los acentos (á é í ó ú ñ ü) y  – — “ ” → .
\usepackage{graphicx}
\usepackage{listings}
\usepackage{hyperref}
\usepackage{xcolor}

% ---- Estilo para código VHDL (apéndice y fragmentos) ----
\lstdefinestyle{vhdl}{
  language=VHDL,
  basicstyle=\ttfamily\scriptsize,
  keywordstyle=\color{black},
  commentstyle=\color{gray}\itshape,
  stringstyle=\color{black},
  showstringspaces=false,
  breaklines=true,
  frame=single,
  rulecolor=\color{black}
}
\lstset{style=vhdl}

\begin{document}

% ============================================================
% TÍTULO Y AUTORES  (los 5 integrantes del equipo — por rellenar)
% ============================================================
\title{Comparación de las funciones de mapeo de la memoria caché\\
mediante simulación VHDL}

\author{Nombre Apellido$^{1}$,%
\and Nombre Apellido$^{1}$,%
\and Nombre Apellido$^{1}$,%
\and Nombre Apellido$^{1}$,%
\and Nombre Apellido$^{1}$}

\authorrunning{Apellido et al.}
\institute{Equipo de Arquitectura de Computadoras,\\
Universidad, Ciudad, País\\
\email{\{correo1, correo2, correo3, correo4, correo5\}@universidad.mx}}

\maketitle

% ============================================================
% RESUMEN  (< 100 palabras — fuente 9pt con sangría, estilo llncs)
% ============================================================
\begin{abstract}
Resumen. Se comparan mediante simulación VHDL-2008 las tres funciones
de mapeo de la memoria caché: directo, asociativo total y asociativo
por conjuntos. Las tres cachés comparten interfaz, memoria principal
y una traza de 20 accesos con un patrón de thrashing, de modo que la
comparación es justa. Las simulaciones se ejecutan con GHDL y las
ondas se inspeccionan con GTKWave. En la traza común la caché directa
alcanza 35\,\% de aciertos, la por conjuntos 50\,\% y la asociativa
total 65\,\%, resultados coherentes con el coste y la complejidad de
cada mapeo.
\end{abstract}

\noindent\textbf{Palabras Clave:} Memoria caché, Mapeo directo,
Asociativo por conjuntos, VHDL, GHDL, Simulación

% ============================================================
% 1. INTRODUCCIÓN  (1–2 cuartillas)
% ============================================================
\section{Introducción}
\label{sec:intro}

La memoria principal es mucho más lenta que el procesador: cerrar una
brecha de este tipo es el motivo de existir de la jerarquía de
memoria, y la caché es su primer y más efectivo nivel
\cite{stallings2019,hennessy2019,tanenbaum2013}. Una caché aprovecha
los principios de localidad temporal y espacial \cite{wikipedia-locality}
para guardar en un almacenamiento
rápido y pequeño una copia de los bloques que el procesador usa con
más frecuencia \cite{patterson2017,bryant2015}.

Cuando se diseña una caché hay que decidir, para cada bloque de la
memoria principal, qué línea de la caché puede alojarlo. Esa decisión
es la \emph{función de mapeo}, y existen tres básicos: mapeo directo
(\emph{direct mapping}), mapeo asociativo total (\emph{associative
mapping}) y mapeo asociativo por conjuntos (\emph{set-associative
mapping})
\cite{stallings2019,wikipedia-cache,null2024}. Cada uno compromete
distinto hardware de comparación y distinta flexibilidad de
reemplazo, de modo que ninguna opción es \emph{a priori} mejor que
las otras: dependen del patrón de accesos.

Este trabajo construye y compara las tres funciones de mapeo en el
mismo entorno de simulación. Las tres cachés están descritas en
VHDL-2008~\cite{ieee1076}, se simulan con GHDL~\cite{ghdl-docs} y sus ondas
se
inspeccionan con GTKWave~\cite{gtkwave-docs}. Para que la comparación
sea justa, los tres módulos exponen \emph{la misma interfaz}, leen de
\emph{la misma memoria principal de comportamiento} y ejecutan
\emph{la misma traza de accesos}, que incluye un patrón de
\emph{thrashing} típico del mapeo directo. El objetivo es doble:
(i) comprobar por simulación el comportamiento teórico de cada
mapeo, y (ii) cuantificar sus diferencias en aciertos y fallos sobre
una carga común.

Desde el punto de vista didáctico conviene aclarar la terminología. El
equipo denomina a los tres métodos ``directo'', ``asociativo'' y
``asociativo por conjuntos'', siguiendo la descripción de funcionamiento
del laboratorio (búsqueda lineal de instrucciones, uso de etiquetas
con comparación en paralelo y división en conjuntos); en la
bibliografía se usa la terminología estándar de Stallings, que es la
que se emplea en el resto del artículo \cite{stallings2019}.

El artículo se estructura como sigue. La sección~\ref{sec:diseno}
describe las tres funciones de mapeo y los módulos VHDL. La
sección~\ref{sec:propuesta} presenta la metodología experimental. La
sección~\ref{sec:simulaciones} muestra las capturas de las
simulaciones, la sección~\ref{sec:resultados} analiza los resultados
cuantitativos y la sección~\ref{sec:conclusiones} concluye y señala
trabajos futuros.

% ============================================================
% 2. DISEÑO TEÓRICO  (3–5 cuartillas)
% ============================================================
\section{Diseño Teórico}
\label{sec:diseno}

\subsection{Contexto: memoria caché y bloques}

Una caché organiza la memoria principal en \emph{bloques} (o
líneas) de varios bytes y guarda esos bloques en sus \emph{líneas}.
En este trabajo se reproducen a escala reducida los parámetros del
Ejemplo 5.1 del capítulo 5 de Stallings \cite{stallings2019} (caché de
64\,kB, memoria principal de 16\,MB y bloques de 4 bytes):
direccionamiento a nivel de byte
con direcciones de 10 bits, palabra y bloque de 4 bytes (por tanto 2
bits de offset), memoria principal de 1\,KiB (256 bloques) y caché de
64\,B (16 líneas). La Tabla~\ref{tab:par} recoge los parámetros y el
reparto de bits de la dirección de cada mapeo.

\begin{table}[t]
\centering
\caption{Parámetros comunes y reparto de la dirección (10 bits).}
\label{tab:par}
\small
\begin{tabular}{llll}
\hline\noalign{\smallskip}
Parámetro & Valor & Mapeo & Reparto de bits\\
\noalign{\smallskip}\hline\noalign{\smallskip}
Dirección & 10 bits (byte) & Directo & tag(4) $\vert$ índice(4) $\vert$ offset(2)\\
Palabra/bloque & 4 bytes & Asociativo & tag(8) $\vert$ offset(2)\\
Mem. principal & 1\,KiB (256 bloques) & Conjuntos & tag(5) $\vert$ conjunto(3) $\vert$ offset(2)\\
Caché & 64\,B (16 líneas) & & \\
\noalign{\smallskip}\hline
\end{tabular}
\end{table}

\subsection{Mapeo directo}

En el mapeo directo cada bloque solo puede ocupar una única línea de
la caché: la línea $i$ del bloque se calcula como $i = j \bmod m$,
donde $j$ es el número de bloque en memoria y $m$ el número de líneas
de la caché \cite{stallings2019,wikipedia-cpu-cache}. En nuestro caso $m=16$, así
que el índice son los 4 bits intermedios de la dirección y los 4 bits
superiores quedan como etiqueta (\emph{tag}) que debe guardarse y
compararse. La búsqueda consiste en un único comparador de 4 bits
sobre la línea seleccionada por el índice, por lo que es el mapeo más
barato y rápido: no hay que recorrer la caché ni decidir víctima.

Su inconveniente es estructural: dos bloques con el mismo índice
interfieren siempre, y si un programa alterna accesos a dos bloques
que comparten índice, cada acceso es un fallo: es el fenómeno de
\emph{thrashing}
\cite{stallings2019,wikipedia-cache,wikipedia-metrics}. La traza de
este trabajo está diseñada para provocarlo con las palabras 0, 16 y
32, que comparten índice.

\subsection{Mapeo asociativo total}

En el mapeo asociativo total cualquier bloque puede ocupar
cualquiera de las líneas \cite{stallings2019}. La dirección se
divide solo en etiqueta (8 bits) y offset, y la caché guarda la
etiqueta de cada línea: en el ciclo de comparación se comparan
\emph{en paralelo} las 16 etiquetas con la solicitada. Esto requiere
16 comparadores de 8 bits y una red de selección (la estructura
clásica es una memoria de contenido o \emph{CAM}, \emph{content
addressable memory}) \cite{stallings2019,hennessy2019}. El
inconveniente es el coste y la potencia de los comparadores en
paralelo; la ventaja es que un bloque nunca expulsa a otro por
colisión de índice: solo sale de la caché cuando hace falta espacio y
se elige una víctima con una política de reemplazo, aquí LRU
(\emph{least recently used}) \cite{stallings2019}.

\subsection{Mapeo asociativo por conjuntos}

El mapeo por conjuntos es el punto intermedio \cite{stallings2019}.
La caché se divide en conjuntos de $k$ líneas (\emph{$k$-way
set-associative}), y un bloque solo puede ocupar cualquiera de las
$k$ líneas de \emph{su} conjunto; dentro del conjunto la elección es
libre. Con $k=1$ se obtiene el mapeo directo y con $k=m$ el asociativo
total \cite{stallings2019,patterson2017}. En nuestro caso hay 8
conjuntos de 2 vías, de modo que la dirección se reparte en etiqueta
(5 bits), conjunto (3 bits) y offset, y solo hacen falta 2
comparadores de 5 bits por acceso. El reemplazo es LRU con 1 bit por
conjunto \cite{wikipedia-cpu-cache}: se recuerda cuál de las dos vías se usó
más recientemente y la víctima es la otra.

\subsection{Reemplazo LRU}

Las cachés asociativas necesitan una política de reemplazo. Se ha
elegido LRU, que descarta el bloque hace más tiempo sin usar y que
aprovecha mejor la localidad temporal que políticas aleatorias o FIFO
\cite{stallings2019,hennessy2019,wikipedia-replacement}. La caché totalmente
asociativa
implementa LRU exacto con un contador de antigüedad por línea (0 =
más reciente), y la por conjuntos implementa LRU exacto con 1 bit por
conjunto al tener solo dos vías. En un diseño real el LRU exacto de
16 vías no es barato y se sustituye por aproximaciones
\cite{hennessy2019}; aquí se justifica por tratarse de un modelo de
simulación de tamaño reducido.

\subsection{Decisiones de diseño y módulos VHDL}

Las decisiones de diseño, todas ellas dirigidas a que la comparación
sea justa, son:
\begin{itemize}
  \item \textbf{Misma interfaz.} Los tres cachés exponen los mismos
    puertos (\texttt{clk}, \texttt{rst}, \texttt{addr}, \texttt{req},
    \texttt{ready}, \texttt{data}, \texttt{hit}, \texttt{estado} y la
    interface de memoria \texttt{mem\_addr}, \texttt{mem\_req},
    \texttt{mem\_data}, \texttt{mem\_ready}), de modo que el banco de
    pruebas y la memoria son intercambiables.
  \item \textbf{Solo lectura.} Se modela el \emph{fetch} de
    instrucciones: no hay bit sucio (\emph{dirty}) ni política de
    escritura \cite{wikipedia-cache-computing}, que queda como trabajo
    futuro (sección~\ref{sec:conclusiones}).
  \item \textbf{Máquina de estados común.}
    $\texttt{IDLE} \rightarrow \texttt{COMPARE} \rightarrow
    (\texttt{FETCH\_MEM} \rightarrow \texttt{DONE}) \mid
    \texttt{DONE}$: se compara, y solo en caso de fallo se pide el
    bloque a la memoria principal, que tiene una latencia de 1 ciclo
    (modelo de comportamiento).
  \item \textbf{Misma memoria principal.} RAM de comportamiento de
    1\,KiB inicializable con un patrón determinista, compartida por
    los bancos de prueba.
\end{itemize}

Los módulos de \texttt{vhdl/mapeo-memoria/src/} son
\texttt{memoria\_principal} (RAM de 256 palabras de 32 bits con
latencia de 1 ciclo y puerto de escritura para los bancos de prueba),
\texttt{cache\_directo} (1 comparador de 4 bits, sin reemplazo),
\texttt{cache\_asociativa} (16 comparadores de 8 bits en paralelo,
contadores LRU) y \texttt{cache\_conjuntos} (2 comparadores de 5
bits, LRU de 1 bit por conjunto). El código completo está en el
apéndice.

% ============================================================
% 3. PROPUESTA EXPERIMENTAL  (1 hoja)
% ============================================================
\section{Propuesta Experimental}
\label{sec:propuesta}

Metodología: herramientas de simulación (GHDL~\cite{ghdl-docs}, sus
fases de análisis, elaboración y ejecución~\cite{ghdl-invoking} con
volcado de ondas~\cite{ghdl-simulation}, GTKWave~\cite{gtkwave-docs}),
estándar VHDL-2008, bancos
de prueba, variables y procedimiento de evaluación. (1 hoja.)

% ============================================================
% 4. SIMULACIONES DE VHDL  (capturas de los módulos más importantes)
% ============================================================
\section{Simulaciones de VHDL}
\label{sec:simulaciones}

Capturas de los módulos más importantes (máximo 2 hojas). Exporta las ondas
de \texttt{vhdl/mapeo-memoria/salidas/*.vcd} a PNG en \texttt{figuras/}.

% Descomenta el bloque cuando exista la captura en figuras/
% \begin{figure}[t]
%   \centering
%   \includegraphics[width=0.9\textwidth]{figuras/captura_modulo.png}
%   \caption{Simulación de <módulo>: descripción breve de lo que se observa.}
%   \label{fig:captura}
% \end{figure}

% ============================================================
% 5. RESULTADOS  (3–5 cuartillas)
% ============================================================
\section{Resultados}
\label{sec:resultados}

Análisis de las simulaciones: tablas comparativas, tiempos, comportamiento
observado, cumplimiento de las pruebas (assert). Interpretación crítica.
(3--5 cuartillas.)

% ============================================================
% 6. CONCLUSIONES  (1 cuartilla)
% ============================================================
\section{Conclusiones}
\label{sec:conclusiones}

Cierre: logros frente a los objetivos, limitaciones y trabajo futuro.
(1 cuartilla.)

% ============================================================
% REFERENCIAS  (>= 15 — indicar exactamente dónde aparece el texto)
% ============================================================
% IMPORTANTE: cada referencia debe permitir ubicar exactamente dónde
% aparece el texto sugerido (capítulo/§/página, o URL + fecha de consulta).
% Anota la ubicación en el campo note de la entrada.
\bibliographystyle{splncs04}
\bibliography{referencias}

% ============================================================
% APÉNDICE — CÓDIGO VHDL
% ============================================================
\appendix
\section{Código VHDL}
\label{app:codigo}

Código completo de los módulos de la unidad (copiado de
\texttt{vhdl/mapeo-memoria/src/}). Usa \texttt{lstlisting}:

% Sustituye por el código real de cada módulo de la unidad.
% \begin{lstlisting}[caption={sumador.vhd --- <descripción>}, label=lst:sumador]
% library ieee;
% use ieee.std_logic_1164.all;
% use ieee.numeric_std.all;
% ...
% \end{lstlisting}

\end{document}
